\chapter{Project Objectives}\label{ch:objectives}

\pagestyle{fancy}

\section{Theme of the Thesis}

As discussed in Chapter~\ref{ch:introduction}, modern healthcare organizations store huge volumes of EHR data that can either not be analyzed at all or, if it is, is rather difficult for interpretation for the majority of people, such as physicians, hospital managers or even patients. This thesis addresses this issue and aims to \emph{develop and validate an analytics solution called HospitAI that converts EHR data into meaningful clinical insights}.

Such an analytics framework, although not yet realized in literature, could build on technologies that currently operate independently from each other. The clustering of \emph{mixed-type} clinical data can be performed using algorithms like K-Prototypes, while the temporal aspect of it, whose clinical significance cannot be represented in a static fashion, is achievable through sequential pattern mining. What is missing today, and is discussed in detail in Chapter~\ref{ch:research}, is a framework that can combine all these pieces. In particular, it should cluster patients into phenotypes, mine trajectories of care in separate clusters, and interpret the mined patterns automatically. The last step is necessary, since the outputs of algorithms used in the previous steps, namely prototype vectors, support counts and event sequences, are completely opaque to physicians and healthcare decision makers, which causes a \emph{semantic gap} between the data mining phase and the target user group. Stratification, per-cluster trajectory mining, and automatic interpretation are the three key concepts of the proposed framework. Since experimentation with those concepts is not possible with real-world data due to legal restrictions, the proposed framework uses the \emph{Synthea} synthetic dataset in its design and evaluation process, while the costly data analysis is performed offline in batch mode, making it possible to focus on the analytics pipeline rather than production-grade application infrastructure. The rest of this chapter formally sets out the objective of this thesis and six subobjectives that form the underlying structure of this work.

\section{Main Objective}

The primary goal of the present work is to \textbf{design, implement, and evaluate the HospitAI platform}, which is an end-to-end framework that transforms raw EHRs into practical clinical knowledge. This primary goal will be reached by creating a software pipeline that chains clustering, sequence mining, and language-model interpretation. The pipeline ingests and processes data (synthetically generated patient data) using unsupervised clustering algorithms, which allows the creation of phenotypes based on complex medical backgrounds. Then, sequential pattern mining algorithms will be applied for care path and treatment trajectory identification. To connect the purely statistical processing mentioned above with clinical decision-making, the platform uses Large Language Models (LLMs) integrated into the constrained interpretation pipelines for understanding the output. Success is defined as the system's ability to autonomously process raw EHR data, perform phenotyping, and visualize care pathways via a user interface.

\section{Secondary Objectives}

The primary goal discussed above is reached through the accomplishment of the following secondary goals, corresponding to the stages of development of the platform and further chapters of this thesis.

\begin{itemize}
    \item[\textbf{I.}] \textbf{Study the state-of-the-art literature and existing algorithms.} Perform an in-depth study of existing methods of clustering, sequence mining and LLMs, to theoretically justify the chosen architecture of the platform. The literature survey should emphasize the limits of the use of individual approaches in processing complex healthcare datasets, revealing the gap which is covered by the current work. Success is measured by the critical comparison of each work studied in the literature survey to the proposed platform, provided in Chapter~\ref{ch:research}.

    \item[\textbf{II.}] \textbf{Develop an integrated, functional prototype of the platform.} Consolidate the analytical building blocks into a single runnable system rather than disconnected scripts. The prototype couples a web-based, role-aware interface for the Manager, the Doctor, and the Patient with a backend application and a relational database, whose data foundation is established in the following objective. Success is measured by a running end-to-end system that loads the dataset, executes the full pipeline on demand, and presents the resulting phenotypes and care pathways.

    \item[\textbf{III.}] \textbf{Design the end-to-end data flow (ETL pipeline) and database architecture.} Extract, transform and load the synthetic raw data into a normalized relational database to filter the noise, establish the relational links between the entities, and generate patient-level clinical features used by the downstream models. This will create a standardized and reproducible data environment that decouples machine learning tasks from the format of the raw data, forming the data backbone of the functional prototype defined above.

    \item[\textbf{IV.}] \textbf{Implement the proposed methodology for patient stratification and clinical trajectory extraction.} Use unsupervised learning techniques to generate homogeneous phenotypes from patients using attributes of a mixed type natively, and use a pattern-mining approach to find frequent and clinically meaningful sequences of events in the history of each individual cluster. Given the generated phenotypes, the combination of the two different approaches will produce one analytical pipeline where the output of the clustering part will feed into the input of the mining part.

    \item[\textbf{V.}] \textbf{Develop an LLM-based conversational and decision-support interface.} Combine an intelligent agent that has two tasks: secure conversion of natural language questions into read-only SQL queries so that non-expert users could query the database, and creation of a clear and constraint-driven clinical summary for each cluster and pathway found. In both tasks, the language model is intentionally constrained so that it speaks about and queries the available information without ever making up numbers and updating the database. The interface succeeds if a user who lacks technical expertise can get the correct and readable answer backed up by the data.

    \item[\textbf{VI.}] \textbf{Rigorously test, validate, and evaluate the system outcomes.} Evaluate the quality of the architecture design based on the internal validation criteria for the clusters and statistical criteria for the sequential patterns, because there is no way to know the ground truth labels for either type of problem. At the same time, strictly verify the safety and accuracy of the LLM-generated responses to prevent clinical hallucinations, which is a prerequisite for deploying the model on the clinical data. These results are presented in Chapter~\ref{ch:testing}.
\end{itemize}

\vspace{0.5cm}

Having established these objectives, Chapter~\ref{ch:research} reviews the literature relevant to each of these analytical techniques and shows where HospitAI stands in the current state of the art.
